%! Author = Sujal Singh
%! Date = 12/3/23

% Preamble
\documentclass[11pt]{beamer}
\title{Electrical Science Presentation}
\author[Aman, Ashish, Dhruv...]{\(|\) Aman Gupta \(|\) Ashish Chauhan \(|\) Dhruv Grover \(|\) Harsh Raj \(|\)\\
    \(|\) Pranav Bisht \(|\) Sujal Singh \(|\)}
\date[Harsh, Pranav, Sujal]{\textbf{Enrollment Numbers:\\}\texttt{0\{00,00,00,00,00,00\}00000000}}

\usetheme{Madrid}
\usecolortheme{spruce}
\setbeamertemplate{navigation symbols}{}
\setbeamertemplate{frametitle continuation}[from second][]

% Packages
\usepackage{amsmath}
\usepackage{lmodern}
\usepackage{siunitx}

% Document
\begin{document}
    \begin{frame}{Electrical Science}
        \begin{center}
            %! suppress = FileNotFound
            \includegraphics[width=80pt]{logo}
        \end{center}\vspace*{-10pt}
        \maketitle
    \end{frame}

    \begin{frame}[t]{Unit 2}
        \textbf{\large Explain the mathematical expression for voltage and current relationship in star and delta
        connection.}
    \end{frame}

    \begin{frame}[t,allowframebreaks]{Unit 3}
        \textbf{\Large Write the working principle of DC motor. Also explain the classification of different DC
        motors.}\\[20pt]

        \begin{center}
            \underline{\smash{\textbf{Working Principle}}}\\[10pt]%
            \textit{``When a current carrying conductor is placed in a magnetic field, it experiences a mechanical
            force''}
        \end{center}

        The construction of both DC Generator and Motor is same, the only difference is here we provide energy to the
        machine.\ As we excite the armature, the conductors of the armature, placed in magnetic field, experience a
        mechanical force and the motor starts moving.\ The direction of rotation is given by Fleming Left Hand Rule.

        \framebreak

        If we stretch the first finger, second finger and thumb of our left hand to be perpendicular to each other,
        and first finger represents the direction of the magnetic field, the second finger represents the direction
        of the current, then the thumb represents the direction of the force experienced by the current carrying
        conductor.

        \\[10pt]

        \begin{center}
            \underline{\smash{\textbf{Types of DC Motor}}}\\[10pt]%
        \end{center}

        There are three types of motors:

        \begin{enumerate}
            \item Series Motor
            \item Shunt Motor
            \item Compound Motor
            \begin{enumerate}
                \item Long Shunt
                \item Short Shunt
            \end{enumerate}
        \end{enumerate}
    \end{frame}

    \begin{frame}[allowframebreaks]{Unit 4}
        \textbf{\Large What are measuring instruments? Also classify the different categories of the measuring
        instrument.}~\\[10pt]

        Devices that are used to measure the magnitude of an electric current with different objectives are known as
        measuring instruments.\ Some common categories of measuring instruments are:\\[10pt]

        \textbf{Voltage measurement units:}
        \begin{itemize}
            \item \underline{Voltmeter}: Measures electrical potential in volts (\SI{}{\volt}).
            \item \underline{\smash{Oscilloscope}}: Displays voltage waveform overtime.
        \end{itemize}

        \\[10pt]

        \textbf{Current measurement units:}
        \begin{itemize}
            \item \underline{Ammeter}: Measures the flow of electrical current in amperes (\SI{}{\ampere}).
        \end{itemize}

        \framebreak

        \textbf{Resistance measurement units:}
        \begin{itemize}
            \item \underline{Ohmmeter}: Measures electrical resistance in ohms ($\Omega$).
            \item \underline{\smash{Wheatstone Bridge}}: Determines an unknown electrical resistance by balancing two
            legs of a bridge circuit.
        \end{itemize}

        \\[10pt]

        \textbf{Power measurement units:}
        \begin{itemize}
            \item \underline{Watt-meter}: Measures electrical power in watts (\SI{}{\watt}).
            \item \underline{\smash{Power Analyser}}: Analyses electrical power quality.
        \end{itemize}

        \\[10pt]

        \textbf{Impedance measurement units:}
        \begin{itemize}
            \item \underline{\smash{Impedance analyzer}}: Measures the impedance of a circuit.
        \end{itemize}
    \end{frame}
\end{document}